\section{Method}
\label{sec:method}

\subsection{Problem formulation}

Let $X=\{x_t\}_{t=1}^{T}$ denote a multivariate industrial time series, where $x_t \in \R^d$ contains all sensor channels at time step $t$. For a given time index $t$, the model receives a past observation window
\begin{equation}
W_t = \{x_{t-L+1}, \ldots, x_t\}.
\end{equation}

The current state label belongs to the set
\begin{equation}
\mathcal{S} = \{\stopmode, \normal, \degone, \degtwo, \fault\},
\end{equation}
which corresponds to the encoded labels $\{1, 5, 7, 9, 3\}$ in the dataset.

Unlike flat classification, we define a directed state graph
\begin{equation}
\mathcal{G}=(\mathcal{S}, \mathcal{E}),
\end{equation}
where $\mathcal{E}$ contains only valid operational transitions. A practical graph for this task is
\begin{equation}
\mathcal{E} = \{
\stopmode \rightarrow \normal,\,
\normal \rightarrow \degone,\,
\degone \rightarrow \degtwo,\,
\degone \rightarrow \fault,\,
\degtwo \rightarrow \fault,\,
\fault \rightarrow \stopmode,\,
\normal \rightarrow \stopmode
\}.
\end{equation}

The model jointly predicts:
\begin{align}
\hat{y}_t &= f_{\st}(W_t), \\
\hat{z}_t &= f_{\fault}(W_t; H), \\
\hat{b}_t &= f_{\mathrm{time}}(W_t),
\end{align}
where $\hat{y}_t$ is the current state, $\hat{z}_t$ is a short-horizon worsening or fault-entry risk within horizon $H$, and $\hat{b}_t$ is a time-to-fault bucket.

\subsection{Overall architecture}

SGPH-Net consists of four parts:
\begin{enumerate}
\item a sensor-group temporal encoder,
\item a current-state recognition head,
\item a state-graph progression head, and
\item a time-to-fault bucket head.
\end{enumerate}

The goal is to retain a compact backbone while placing the main inductive bias in the output structure.

\subsection{Sensor-group temporal encoder}

The raw channels are divided into physically meaningful groups:
\begin{itemize}
\item speed group,
\item depth group,
\item current group,
\item load and pressure group.
\end{itemize}

Each group is encoded by a lightweight temporal block. In the initial implementation, this block can be instantiated as a temporal convolutional network (TCN) with gated residual layers. Let the group embeddings be
\begin{equation}
h_t^{(g)} = \phi_g(W_t), \quad g \in \{1, \ldots, G\}.
\end{equation}
The final shared representation is
\begin{equation}
h_t = \psi\left(h_t^{(1)}, \ldots, h_t^{(G)}\right),
\end{equation}
where $\psi(\cdot)$ denotes concatenation followed by a fusion multilayer perceptron.

\subsection{Current-state recognition head}

The current state distribution is predicted as
\begin{equation}
p_t = \softmax(W_y h_t + b_y),
\end{equation}
where $p_t \in \R^{|\mathcal{S}|}$.

This head is responsible for five-state recognition. However, unlike a flat classifier, its output is coupled to the progression head described next.

\subsection{State-graph progression head}

We define a learnable transition score matrix $T_t \in \R^{|\mathcal{S}| \times |\mathcal{S}|}$ conditioned on the current window:
\begin{equation}
T_t = \mathrm{reshape}(W_T h_t + b_T).
\end{equation}

To suppress invalid transitions, we introduce a binary adjacency mask $M$ derived from $\mathcal{E}$:
\begin{equation}
\tilde{T}_t = T_t \odot M.
\end{equation}

The one-step progression distribution is then approximated by
\begin{equation}
q_t = p_t \tilde{T}_t.
\end{equation}

The short-horizon worsening score is predicted as
\begin{equation}
\hat{z}_t = \sigmoid(W_h h_t + b_h).
\end{equation}

This design couples current-state representation with transition-aware warning, which is more appropriate than optimizing warning as a detached auxiliary label.

\subsection{Time-to-fault bucket head}

For samples that have not yet entered the final fault state, SGPH-Net predicts a bucketized time-to-fault distribution:
\begin{equation}
\hat{b}_t = \softmax(W_b h_t + b_b).
\end{equation}

In the current draft, we use the following buckets:
\begin{itemize}
\item 0--4 s,
\item 4--12 s,
\item 12--20 s,
\item 20--40 s,
\item no fault within horizon.
\end{itemize}

Bucketization is preferred over direct regression because the dataset is limited and highly imbalanced, especially around the second-level degradation state.

\subsection{Training objective}

The total loss is defined as
\begin{equation}
\loss = \loss_{\st} + \lambda_1 \loss_{\mathrm{prog}} + \lambda_2 \loss_{\mathrm{hazard}} + \lambda_3 \loss_{\mathrm{time}}.
\end{equation}

Here:
\begin{itemize}
\item $\loss_{\st}$ is the class-balanced state-recognition loss,
\item $\loss_{\mathrm{prog}}$ penalizes mass assigned to invalid edges,
\item $\loss_{\mathrm{hazard}}$ supervises short-horizon worsening prediction,
\item $\loss_{\mathrm{time}}$ supervises the time-to-fault bucket prediction.
\end{itemize}

A simple invalid-transition penalty is
\begin{equation}
\loss_{\mathrm{prog}} = \sum_{i,j} (1-M_{ij}) \tilde{T}_{t,ij}^{2}.
\end{equation}

The proposed formulation is intentionally compact. The novelty lies in the state-graph-constrained progression learning view, not in building a heavier backbone.
